\chapter{Weak convergence, continuity sets, and tightness}\label{sec:measure-convergence-tools}
We now supply the compactness argument used throughout GMT and explain when weak convergence controls masses. All measures are nonnegative Radon measures on an open $U\subset\mathbb R^n$ unless otherwise specified. The main source is \cite{Sim18}, Chapter 1, Section 5, especially Theorem 5.15, pp.~38--39. The countable testing argument below uses the Euclidean setting explicitly. Examples \ref{ex:escaping-mass} and \ref{ex:weak-borel-discontinuity} explain why the qualifications are needed.

\begin{lemma}[A countable family of compactly supported tests]\label{lem:countable-radon-tests}
There are compact sets $K_m\subset\operatorname{int}K_{m+1}\subset U$ exhausting $U$, and a countable family $\mathcal D\subset C_c(U)$, such that every $f\in C_c(U)$ supported in $K_m$ is a uniform limit of members of $\mathcal D$ supported in $K_{m+1}$.
\end{lemma}
\begin{proof}
One may take $K_m=\{x\in U:|x|\leq m,\ \operatorname{dist}(x,\mathbb R^n\setminus U)\geq1/m\}$, with distance to the empty set interpreted as infinity. These are compact, lie in the interior of the next set, and their interiors cover $U$. Choose fixed cutoffs $0\leq\eta_m\leq1$, equal to one on $K_m$, with support in $\operatorname{int}K_{m+1}$, using Lemma \ref{lem:radon-compact-cutoffs}.

There is a countable family $\mathcal P$ of compactly supported continuous functions on $\mathbb R^n$ dense for the uniform norm in $C_c(\mathbb R^n)$: use piecewise multilinear interpolation on dyadic grids, with finitely many nonzero rational vertex values. Uniform continuity controls interpolation error on sufficiently small cubes; approximation of the finitely many vertex values by rationals adds arbitrarily small uniform error. A function in $C_c(U)$ extends continuously by zero to $\mathbb R^n$. If its support is in $K_m$, uniform approximants $q_j\in\mathcal P$ give $\eta_mq_j\to\eta_mf=f$, with all supports in $K_{m+1}$. The countable union $\mathcal D=\bigcup_m\eta_m\mathcal P$ has the asserted property. Empty initial compact sets cause no difficulty.
\end{proof}

\begin{lemma}[Diagonal extraction of a positive limit functional]\label{lem:measure-diagonal-functional}
Suppose $\sup_j\mu_j(K)<\infty$ for every compact $K\subset U$. There are a subsequence $\mu_{j_k}$ and a positive linear functional $L:C_c(U)\to\mathbb R$ with
\[
 \int f\,d\mu_{j_k}\longrightarrow L(f)\qquad(f\in C_c(U)).
\]
\end{lemma}
\begin{proof}
Enumerate the family $\mathcal D$ from Lemma \ref{lem:countable-radon-tests}. For each of its functions the integrals form a bounded real sequence, because its support is compact. Successive subsequences and the diagonal construction give simultaneous convergence for every member of $\mathcal D$. For arbitrary $f$ supported in $K_m$, approximate it uniformly by $g\in\mathcal D$ supported in $K_{m+1}$. If $M_m=\sup_j\mu_j(K_{m+1})$, then
\[
 \left|\int f\,d\mu_{j_k}-\int f\,d\mu_{j_\ell}\right|
 \leq2M_m\|f-g\|_\infty+
 \left|\int g\,d\mu_{j_k}-\int g\,d\mu_{j_\ell}\right|.
\]
First make the uniform error small and then take $k,\ell$ large. Thus the integrals of $f$ are Cauchy and define a finite limit $L(f)$. Limits preserve real linear combinations and nonnegativity. This proves the assertion. Sequential extraction here is justified by the countable tests, not merely by a compactness assertion in a possibly nonmetrizable dual space.
\end{proof}

\begin{proposition}[Upper semicontinuity on compact sets]\label{prop:measure-compact-usc}
If $\mu_j\rightharpoonup\mu$ in $U$, then
\[
 \limsup_j\mu_j(K)\leq\mu(K)
\]
for every compact $K\subset U$.
\end{proposition}
\begin{proof}
By outer regularity and local finiteness, choose an open neighborhood $O$ of $K$ with compact closure in $U$ and $\mu(O\setminus K)<\varepsilon$. A cutoff $0\leq\eta\leq1$ supported in $O$ and equal to one on $K$ gives
\[
 \limsup_j\mu_j(K)\leq\lim_j\int\eta\,d\mu_j
 =\int\eta\,d\mu\leq\mu(O)\leq\mu(K)+\varepsilon.
\]
Let $\varepsilon\downarrow0$. Compactness both permits the cutoff and makes all these integrals finite. Compare the open-set inequality in Proposition \ref{prop:measure-lsc}.
\end{proof}

\begin{corollary}[Convergence of masses on continuity sets]\label{cor:measure-continuity-sets}
If $\mu_j\rightharpoonup\mu$ in $U$ and $E\subset U$ is Borel with compact closure in $U$ and $\mu(\partial E)=0$, then $\mu_j(E)\to\mu(E)$. In particular, for a fixed center $x\in U$, this holds for $E=B_r(x)$ at all but countably many radii with $\overline B_r(x)\subset U$.
\end{corollary}
\begin{proof}
Proposition \ref{prop:measure-lsc} and Proposition \ref{prop:measure-compact-usc} give
\[
 \mu(\operatorname{int}E)\leq\liminf_j\mu_j(E)
 \leq\limsup_j\mu_j(E)\leq\mu(\overline E).
\]
The boundary hypothesis makes both endpoints equal to $\mu(E)$. For the ball statement, distinct centered spheres are disjoint. Inside any fixed compact ball of finite measure, only finitely many of them can have mass at least $1/k$ for a fixed positive integer $k$. Taking the union over $k$ and over a countable exhaustion of the admissible radii gives at most countably many exceptional spheres.
\end{proof}

\begin{definition}[Uniform tightness and bounded-continuous testing]\label{def:uniform-tightness}
A sequence of finite Radon measures $\mu_j$ on $U$ is uniformly tight if, for every $\varepsilon>0$, some compact $K\subset U$ satisfies $\sup_j\mu_j(U\setminus K)<\varepsilon$. Convergence against all $f\in C_b(U)$, the bounded continuous functions, is called narrow convergence here. It is stronger than the compact-test convergence in Definition \ref{def:weak-measure}; the two agree on uniformly tight families with uniformly bounded total mass, as proved next.
\end{definition}

\begin{theorem}[Tightness upgrades compact tests to bounded tests]\label{thm:tight-vague-to-narrow}
Suppose $\mu_j\rightharpoonup\mu$ in $U$, $\sup_j\mu_j(U)<\infty$, and the sequence is uniformly tight. Then $\mu$ is finite and $\mu_j$ converge narrowly to $\mu$. In particular $\mu_j(U)\to\mu(U)$; if all $\mu_j$ are probability measures, so is $\mu$. The compact-test convention is that of \cite{Sim18}, Theorem 5.15; the additional tightness hypothesis here controls the tails, by the argument below.
\end{theorem}
\begin{proof}
The open-set lower-semicontinuity inequality gives $\mu(U)\leq\liminf_j\mu_j(U)<\infty$. Given $\varepsilon>0$, tightness supplies a compact set controlling all $\mu_j$ tails; inner regularity of the finite measure $\mu$ supplies another compact set controlling its tail. Their union $K$ controls both. Choose $0\leq\eta\leq1$ compactly supported and equal to one on $K$. For $f\in C_b(U)$, the test $\eta f$ is compactly supported, while
\[
 \left|\int(1-\eta)f\,d\mu_j\right|
 +\left|\int(1-\eta)f\,d\mu\right|\leq2\varepsilon\|f\|_\infty.
\]
Take the limit in the compactly supported part and then let $\varepsilon\downarrow0$. The constant test $f=1$ gives convergence of total masses. In contrast, the escaping measures in Example \ref{ex:escaping-mass} are not uniformly tight.
\end{proof}

\begin{corollary}[Sequential compactness for tight probability measures]\label{cor:tight-probability-compactness}
Every uniformly tight sequence of Radon probability measures on $U$ has a narrowly convergent subsequence whose limit is a Radon probability measure. Conversely, a narrowly convergent sequence of Radon probability measures is uniformly tight.
\end{corollary}
\begin{proof}
The local bounds required by Theorem \ref{thm:radon-compactness} hold with upper bound one. Extract a compact-test limit and apply Theorem \ref{thm:tight-vague-to-narrow}.

Conversely, suppose $\mu_j$ converge narrowly to the probability measure $\mu$. Choose compact $K$ with $\mu(K)>1-\varepsilon$ and a cutoff $0\leq\eta\leq1$ equal to one on $K$, with compact support $L$. For all sufficiently large $j$, $\int\eta\,d\mu_j>1-2\varepsilon$, so $\mu_j(U\setminus L)<2\varepsilon$. Each of the finitely many remaining probability measures has a compact set carrying all but $2\varepsilon$ of its mass by inner regularity. Adding those sets to $L$ gives a single compact set for the whole sequence. Letting the initial $\varepsilon$ be arbitrarily small proves tightness.
\end{proof}

\begin{proposition}[Smooth densities approximate any Radon measure locally]\label{prop:radon-mollification}
Let $\mu$ be a Radon measure on $\mathbb R^n$ and $0\leq\rho\in C_c^\infty(B_1)$ with $\int\rho=1$. For $\varepsilon>0$ set
\[
 f_\varepsilon(x)=\int\varepsilon^{-n}\rho((x-y)/\varepsilon)\,d\mu(y),
 \qquad d\mu_\varepsilon=f_\varepsilon\,dx.
\]
Then $f_\varepsilon$ is nonnegative and smooth, $\mu_\varepsilon$ is Radon, and $\mu_\varepsilon\rightharpoonup\mu$ as $\varepsilon\downarrow0$. Finite total mass of $\mu$ is not required.
\end{proposition}
\begin{proof}
On a fixed compact set of $x$ values, every contributing $y$ lies in a fixed larger compact set. Local finiteness bounds the integrals of every derivative of the smooth kernel there, justifying differentiation under the integral and proving smoothness. For $\varphi\in C_c(\mathbb R^n)$, Fubini and a change of variables give
\[
 \int\varphi\,d\mu_\varepsilon
 =\int\left(\int\varphi(y+\varepsilon z)\rho(z)\,dz\right)d\mu(y).
\]
The inner function converges uniformly to $\varphi(y)$ by uniform continuity. For $0<\varepsilon\leq1$, all these functions have support in the compact one-neighborhood of $\operatorname{spt}\varphi$, whose $\mu$ mass is finite. This proves convergence. Smooth locally bounded densities define Radon measures. The conclusion generalizes the concentration example and the diffuse Cantor approximations in Examples \ref{ex:approximate-identity-dirac} and \ref{ex:cantor-density-approximation}.
\end{proof}

\begin{proposition}[Finite atomic measures also approximate locally]\label{prop:radon-atomic-approximation}
Every Radon measure $\mu$ on $\mathbb R^n$ is a compact-test limit of finite nonnegative atomic measures. Consequently compact-test convergence alone preserves neither absolute continuity nor atomicity.
\end{proposition}
\begin{proof}
Partition the half-open cube $[-j,j)^n$ into finitely many half-open cubes $Q$ of side $1/j$, and choose each cube center $x_Q$. Define $\nu_j=\sum_Q\mu(Q)\delta_{x_Q}$; local finiteness makes this a finite measure. For a test $\varphi$ supported in a compact $K$, sufficiently large outer cubes contain the compact one-neighborhood of $K$. Cubes contributing either to $\varphi(x_Q)$ or to its $\mu$ integral lie in that neighborhood once $\sqrt n/j<1$. Thus
\[
 \left|\int\varphi\,d\nu_j-\int\varphi\,d\mu\right|
 \leq\mu(K_1)\,\omega_\varphi(\sqrt n/j),
\]
where $K_1=\{x:\operatorname{dist}(x,K)\leq1\}$ and $\omega_\varphi$ is a modulus of continuity. This tends to zero. Taking Lebesgue measure as the limit gives non-atomic limits of atomic measures; Proposition \ref{prop:radon-mollification} gives smooth-density approximations of singular limits.
\end{proof}
